\section{Implementation}
\label{sec:implementation}

We implement \systemname\ as a drop-in runtime layer on top of HuggingFace
Transformers~\cite{wolf2020transformers}, with $\sim$4.5K lines of Python
and $\sim$0.6K lines of CUDA. The same Python codebase compiles for both
NVIDIA (PyTorch 2.3 / CUDA 12.2) and Ascend (CANN 7.0): a thin $\sim$150-LoC
device shim hides the difference between CUDA streams and the Ascend runtime
equivalents, so the controller and forecaster are device-independent. The
implementation modifies neither the model architecture nor the routing/gating
semantics, and can be swapped in front of any MoE model that follows
HuggingFace's \texttt{*MoeSparseMoeBlock} convention.

The forecaster is trained from $(\mathbf{t}, \ell, S, \mathbf{a})$ tuples
logged in a fixed-size pinned ring buffer (default 64K entries) during
serving. Retraining runs on the controller's CPU thread once per epoch
(10K requests or 60 seconds, whichever comes first) and never touches the
GPU. On a previously unseen model the buffer is empty for the first epoch
and the forecaster falls back to the model's own pregate top-$k$, so
\systemname\ degrades to a baseline-equivalent state at cold start rather
than failing.

\niparagraph{Latency attribution.}
For each MoE layer the runtime emits four \texttt{torch.cuda.Event}
timestamps---compute begin, compute end, first stall on a missing expert,
and stall release---and aggregates them in a per-request ring buffer. The
difference between stall release and stall begin is the \emph{exposed
waiting} component; the time spent inside the swap-in path is the
\emph{cache-miss} component; compute time is reported separately. The
breakdown plotted in \autoref{sec:evaluation}, including
Tables~\ref{tab:Dynamic_baseline}--\ref{tab:Dynamic_CCR}, is read directly
from these counters. Seeds are fixed for ShareGPT bucketing and forecaster
training; we will release the runtime, device shim, training pipeline, and
benchmark harness at the camera-ready URL.
